\documentclass[12pt]{article}
\usepackage{amsmath} % flere matematikkommandoer
\usepackage[utf8]{inputenc} % æøå
\usepackage[T1]{fontenc} % mere æøå
\usepackage[danish]{babel} % orddeling
\usepackage{verbatim} % så man kan skrive ren tekst
\usepackage[all]{xy} % den sidste (avancerede) formel i dokumentet
\title{En kort \LaTeX-skabelon}
\author{Henrik og Henrik}
\begin{document}
\maketitle
\begin{abstract}
Dette dokument er en \LaTeX-skabelon, hvor det kan betale sig at studere tex-koden
og pdf-dokumentet sammen! Vi har bestræbt os på at holde skabelonen kort; se \
cite{n} for en meget mere grundig introduktion.
\end{abstract}
\tableofcontents
\section{Indledning}
Et \LaTeX-dokument starter med \verb!\documentclass!, derefter en {\em preample},
der fastlægger generelle egenskaber ved dokumentet, så \verb!\begin{document}!
efterfulgt af det egentlige dokument og til sidst \verb!\end{document}!.
\section{Lidt om matricer}
I \LaTeX kan man skrive en matrix vha \verb!array! som
$$
\mathbf A=
\left(\begin{array}{cc}
1&7\\9&13
\end{array}\right).

\begin{pmatrix}
1&7\\9&13
\end{pmatrix}
\quad
\begin{bmatrix}
1&7\\9&13
\end{bmatrix}.
$$
Hvis pakken \verb!inputenc! inkluderes kan man skrive æ, ø, å direkte fra et dansk
tastatur.
\section{Almindelig tekst}
Bemærk, at flere mellemrum mellem ord eller bogstaver er det samme som 1 blanktegn,
men en tom linje giver en ny paragraf.
Specialtegn skrives i \LaTeX\ vha ``backslash'', fx giver \verb!\$! et dollartegn
(\$). Efter \verb!%! kan skrives kommentarer på resten af linjen.
\section{Matematik}
Man skifter mellem almindelig tekst og matematiktekst med dollartegn, fx skrives
ligningen $Ax=b$ vha \verb!$Ax=b$! og det er pænere end Ax=b, ikke sandt?\ Lidt
mere blæret er vel $\mathbf {Ax}=\mathbf{b}$ når vi har med matricer og vektorer at
gøre.
Formler på egne linjer skrives som matematikkode mellem \verb!$$! og \verb!$$!,
(eller ækvivalent \verb!\begin{equation*}! og \verb!\end{equation*}!) som fx
$$
\begin{bmatrix}
1&7\\9&13
\end{bmatrix}\begin{bmatrix}x_1\\x_2\end{bmatrix}=\begin{bmatrix}3\\4\end{bmatrix}.
$$
For ``pæne'' formler over flere linjer kan \verb*!\begin{align*}...\end{align*}!
bruges sammen med \verb!&! (alignment) og \verb!\\! (linjeskift), som fx
\begin{align*}
2x_1+3x_2-x_3&=1\\4x_1+7x_2+x_3&=3\\
7x_1+10x_2-4x_3&=4
\end{align*}
Sæærlige bogstavsnit osv: Vi har allerede brugt \verb!\mathbf!, men nogle gange
skrives $\mathcal U$ for et underrum i stedet for bare $U$.
\section{Referencer til formler osv}
Man kan referere tiltidligere (eller senere) formler i sit dokument ved at bruge \
verb!\label{...}! og \verb!\ref{...}!, fx kan man giver ligningen nedenfor en
etiket med fx \verb!\label{henriks-ligning}! (men så skal man huske at erstatte
dollartegnene med \verb!\begin{equation}! og \verb!\end{equation}!):
\begin{equation}
\label{henriks-ligning}
\begin{bmatrix}
1&7\\9&13
\end{bmatrix}\begin{bmatrix}x_1\\x_2\end{bmatrix}=\begin{bmatrix}3\\4\end{bmatrix}.
\end{equation}
Når ligningen har fået en etiket, kan man henvise til den ved at skrive \verb!\
ref{henriks-ligning}! som giver \ref{henriks-ligning}. Når men refererer til en
ligning kan man med fordel bruge \verb!\eqref{henriks-ligning}!, så kommer der
automatisk paranteser omkring 1, fx \eqref{henriks-ligning}.
Bemærk, at det også er muligt at lave sætninger (Theorems), som man kan referere
til. Det fører for vidt at beskrive det her. Se \cite[Section 3.9]{n}.
\section{Hvad er der mere brug for?}
\begin{itemize}
\item Kvadrat skrives som $2^2$ og kvadratrod skrives $\sqrt{2}$
\item Prikproduktet skrives som fx $x\cdot y$, eller $\mathbf x\cdot\mathbf y$.
\item Numerisk værdi og længde som $|x|$ og $||\mathbf x||$
\item Man kan skrive tekst, der ikke oversættes, men videregives som ren tekst. For
enkelte ord kan man bruge \verb!\verb!, mens man for længere afsnit (over flere
linjer) bør bruge \verb!verbatim!
\begin{verbatim}
tekst over
flere linjer
som kan
rykkes ind
\end{verbatim}
\end{itemize}
\section{Exercise 1.1.5}
Her følger et eksempel på indskrivning af en opgave fra bogen. Opgaven går ud på at
bestemme pivotvariable og frie variable for systemet
\begin{align*}
x_1-2x_2-10x_4&=0\\
-x_2+x_3-2x_4&=1\\
2x_3+4x_4&=6
\end{align*}
og dernæst at løse det. Systemet omskrives til matrixform:
$$
\left[\begin{array}{rrrr|r}
1&-2&0&-10&0\\
0&-1&1&-2&1\\
0&0&2&4&6
\end{array}\right]
$$
Pivotelementerne står i positionerne $(1,1)$, $(2,2)$ og $(3,3)$. Dermed er $x_4$
en fri variabel. Sætter vi $x_4=t$ finder vi ved ``back substitution'': $2x_3+4t=6$
dvs
$$
x_3=\frac{6}{2}-\frac{4t}{2}=3-2t.
$$
Dette giver (svarende til anden ligning):
$$
-x_2+x_3-2x_4=1 \quad \text{dvs}\quad x_2=1-\left(3-2t\right)+2t=-2+4t,
$$
og endeligt
$$
x_1=2x_2+10x_4=2\left(-2+4t\right)+10t=-4+18t.
$$
Dette kan skrives som
$$
\begin{bmatrix}x_1\\x_2\\x_3\\x_4\end{bmatrix}=\begin{bmatrix}-4+18t\\-2+4t\\3-2t\\
t\end{bmatrix}
=\begin{bmatrix}-4\\-2\\3\\0\end{bmatrix}+t\begin{bmatrix}18\\4\\-2\\1\end{bmatrix}
$$
\section{Rækkeoperationer}
Neden for vises to måder, hvorpå man kan vise rækkeoperationer. Det simpleste er
nok
som fx:
Ved at bruge rækkeoperationerne $4r_1+r_2 \to r_2$ og $-r_1+r_3 \to r_3$ får man
$$
\begin{bmatrix}
1 & 2 & -1 & 4 & 5 \\
-4 & 2 & 6 & 0 & -2 \\
1 & 3 & 0 & 4 & 10
\end{bmatrix}
\to
\begin{bmatrix}
1 & 2 & -1 & 4 & 5 \\
0 & 10 & 2 & 16 & 18 \\
0 & 1 & 1 & 0 & 5
\end{bmatrix}.
$$
Dette kan sevlfølgeligt typesættes mere avanceret, som fx:
Ved at bruge rækkeoperationer får man
$$
\begin{array}{rcl}
{
\mathbf{M} =
\left(\!\!\!\!
\begin{array}{rrrr|r}
1 & 2 & -1 & 4 & 5 \\
-4 & 2 & 6 & 0 & -2 \\
1 & 3 & 0 & 4 & 10
\end{array}
\!\!\right)
}
&
\xymatrix@C=15ex{
\ar[r]^-{\small
\begin{array}{r}
4\mathbf{r}_1+\mathbf{r}_2 \to \mathbf{r}_2 \\
-\mathbf{r}_1+\mathbf{r}_3 \to \mathbf{r}_3
\end{array}
} &
}
&
{
\left(\!\!
\begin{array}{rrrr|r}
1 & 2 & -1 & 4 & 5 \\
0 & 10 & 2 & 16 & 18 \\
0 & 1 & 1 & 0 & 5
\end{array}
\!\!\right).
}
\end{array}
$$
\begin{thebibliography}{99} % de 99 laver blot lidt plads til venstre for
referencenumeret
\bibitem{n} Tobias Oetiker, Hubert Partl, Irene Hyna and Elisabeth Schlegl, The Not
So Short Introduction to \LaTeXe, Version 5.01, April 06, 2011
\end{thebibliography}
\end{document}
